\documentclass[phyai.tex]{subfiles}
\begin{document}
\chapter*{Preface}\label{chap:motivation} % (fold)
% chapter Motivation (end)
\indent OK, let's get started. What is Artificial Intelligence? I do not know; and even after writing this book, I still do not know. Is it a super-machine that can answer all our questions? Or is it just a statistical machine, sometime quite silly? But let's clarify a bit; how do we define "intelligence"? And if machines can learn the way humans learn (through patterns and constantly updating neural weights), can machines think?
\\ \\
\textit{"Things change, and we can't do nothing about it. Just let it flow, being lighter. Overflowing."}
\\ \\
\indent Third year. It is now the summer of second semester of university. This is the time when people usually think about their future careers, my friends around me are all preparing for internships somewhere and starting jobs. And what about me? Perhaps I am dreaming, and I do not want to wake up yet. I am too ignorant to be a theorist (even in my favorite fields like Physics or Mathematics), but too dreamy to be a computer engineer (which will be printed on my degree). I am stuck between two worlds, and I do not belong to either.
\begin{center}
	\textit{"This might be a good thinking game."}
\end{center}
I remember that day. It was a beautiful sunny day, and I was in the university library. Oh, the next test was in two hours, but how could I cram the enormous amount of knowledge taught in two months into just one hour (not to mention the bus waiting time)? My mind was filled with Fujii Kaze's songs; and yeah, another failed test would not a huge problem, after all. I let my mood drift; and then accidentally clicked on a video of Geoffrey Hinton's banquet speech in Nobel Prize award ceremony.
\begin{center}
	\textit{"They are no longer science fiction."}
\end{center}
This is the first time in my life I have learned that the Artificial Intelligence was initially inspired by Physics. I was incredibly surprised, and felt a little proud and happy. I have finally found the first man on Earth, who truly cared and dedicated his life to pursue the question of how brain works. The coolest about him is that he applied many tools from Physics to describe his ideas in Biology and Artificial Intelligence. Perhaps if we understand science deeply enough, the boundaries between branches will almost fade and disappear.
\\ \\
\indent Back to reality, my life is chaotic and oscillate like a sine wave. If you expect me to be like a superstar, a typical successful figure at university, you will be very disappointed. But I truly love "thinking" from the bottom of my heart. The spirit of thinking and curiosity about everything has driven me to live to this day, to type these lines to you.
\begin{center}
	\textit{{"I am never tired of learning new things."}}
\end{center}
Currently, we are at the first step in understanding what Artificial Intelligence from a physical perspective. I wrote this book as a very-easy-understand to Statistical Physics (in my own way); and our journey is still very long. Perhaps there is no destination, but that is precisely the meaning of scientific research.
\begin{center}
	\textit{"This must be the best thinking game ever."}
\end{center}
My plan is quite simple. Firstly, this book is written to cover the basic concepts and terminology in Statistical Physics. There is no Artificial Intelligence in this book, but there are many connections between them, and they share a common origin such as equilibrium, free energy or Ising model. To fully enjoy this thinking game, I recommend reading my book with just paper, pen, calculator and book (this is exactly how I learn everything too). The best way is following each line that I have calculated and proven, to build the most important thing that humans need: "intuition". Secondly, after completing this book, I plan to start an online blog. The content will cover around how to bridge the gap between Physics and Artificial Intelligence. I will primarily base my work on the Hopfield Network and Boltzmann Machine; or furthermore, any Energy-Based Models (EBM).
\\ \\
\indent But I have not decided which platform to use yet, so the best way to keep track of my work is to follow me on Facebook. I will be constantly updating new announcements about my work.
\begin{flushleft}
	My Facebook's address: \href{https://www.facebook.com/profile.php?id=61577979788883}{Link}
\end{flushleft}
\begin{flushleft}
	My Email address: \url{tinvu1309@gmail.com}
\end{flushleft}
If you are interested and want to meet me in person, feel free to contact me and we could meet somewhere near my university: \textit{University of Engineering and Technology, Vietnam National University (Hanoi).}

\end{document}


